\title{Chain-of-Thought Prompting Elicits Reasoning in Large Language Models}

\begin{abstract}
We investigate how producing a \textit{chain of thought}—a sequence of intermediate reasoning steps—greatly enhances the capacity of large language models to tackle complex reasoning problems. Specifically, we demonstrate that such reasoning skills naturally arise in sufficiently large language models through a straightforward technique called \textit{chain-of-thought prompting}, which involves including a few chain-of-thought examples as exemplars during prompting.

Evaluations on three large language models reveal that chain-of-thought prompting boosts performance across a variety of arithmetic, commonsense, and symbolic reasoning tasks. The observed improvements are substantial. For example, prompting a PaLM 540B model with only eight chain-of-thought examples attains state-of-the-art results on the GSM8K benchmark for math word problems, outperforming even finetuned GPT-3 models equipped with a verifier.
\end{abstract}

\section{Introduction}
Recent advances in NLP have been driven by language models \citep[][\textit{inter alia}]{peters-etal-2018-deep,devlin-etal-2019-bert,brown2020language}. Increasing the size of these models has been shown to bring various advantages, including enhanced accuracy and better sample efficiency \citep[\textit{inter alia}]{kaplan2020scaling,brown2020language}. Nevertheless, merely enlarging model size has not been adequate to achieve top performance on difficult tasks involving arithmetic, commonsense, or symbolic reasoning \citep{rae2021scaling}.

In this paper, we investigate how large language models' reasoning capabilities can be unlocked using a simple approach inspired by two key ideas. First, methods for arithmetic reasoning benefit from producing natural language explanations that lead to the final solution. Previous research has enabled models to generate natural language intermediate steps either by training from scratch \citep{ling-etal-2017-program} or by fine-tuning pretrained models \citep{cobbe2021training}, as well as through neuro-symbolic approaches that utilize formal languages in lieu of natural language \citep{roy-roth-2015-solving, chiang-chen-2019-semantically,amini-etal-2019-mathqa, chen2019neural}. Second, large language models present an exciting opportunity for in-context few-shot learning through \emph{prompting}. That is, instead of fine-tuning a distinct language model checkpoint for each new task, one can simply “prompt” the model with a few input-output pairs that illustrate the task. Notably, this strategy has achieved success across various straightforward question-answering tasks \citep{brown2020language}.

Nonetheless, both of these approaches have significant drawbacks. For rationale-augmented training and finetuning techniques, generating a substantial collection of high-quality rationales is expensive and far more complex than producing simple input--output pairs typically used in conventional machine learning. The traditional few-shot prompting strategy employed in \citet{brown2020language} tends to perform poorly on tasks demanding reasoning skills and often shows limited improvement even as the scale of the language model increases \citep{rae2021scaling}. In this study, we integrate the advantages of these two approaches while circumventing their respective limitations. Concretely, we investigate the capability of language models to execute few-shot prompting on reasoning tasks by utilizing prompts composed of triples: $\langle$input, \emph{chain of thought}, output$\rangle$. A \emph{chain of thought} refers to a sequence of intermediate natural language reasoning steps that culminate in the final answer, and we call this method \textit{chain-of-thought prompting}. An example of such a prompt is depicted in \cref{fig:pull-figure}.

We provide empirical results across benchmarks in arithmetic, commonsense, and symbolic reasoning, demonstrating that chain-of-thought prompting surpasses standard prompting, occasionally by a substantial margin. One such outcome is presented in \cref{fig:pull-bar-chart}—on the GSM8K benchmark for math word problems \citep{cobbe2021training}, chain-of-thought prompting with \palm{} 540B significantly outperforms traditional prompting and sets a new state-of-the-art. A prompting-only approach is valuable because it does not rely on large training datasets and enables a single model checkpoint to handle multiple tasks without compromising generality. This study highlights how large language models can effectively learn from just a few examples accompanied by natural language explanations about the task, contrasting with the conventional method of learning patterns in inputs and outputs through extensive training data.

\section{Chain-of-Thought Prompting} When tackling a complex reasoning problem, such as a multi-step math word problem, individuals typically reflect on their own cognitive process. They often break down the problem into smaller intermediate steps, solving each sequentially before arriving at the final answer: \textit{``After Jane gives 2 flowers to her mom she has 10 $\ldots$ then after she gives 3 to her dad she will have 7 $\ldots$ so the answer is 7.''} This paper aims to equip language models with the capability to produce a comparable \textit{chain of thought}—a logical sequence of intermediate reasoning steps that culminate in the correct answer. We demonstrate that language models of sufficient scale can generate chains of thought when few-shot prompting includes examples demonstrating chain-of-thought reasoning.

\cref{fig:pull-figure} presents an example where a model generates a chain of thought to correctly solve a math word problem that it would have otherwise answered incorrectly. Here, the chain of thought resembles a solution and can be interpreted as such, but we prefer the term chain of thought to emphasize that it simulates a step-by-step reasoning process to reach the conclusion (furthermore, solutions or explanations generally follow the final answer \citep[][\textit{inter alia}]{narang2020wt5,wiegreffe2021reframing,lampinen2022can}).

Chain-of-thought prompting offers several appealing advantages as a method to enhance reasoning capabilities in language models.

\begin{enumerate}[topsep=1pt,itemsep=0ex]
    \item Firstly, chain of thought conceptually enables models to break down multi-step problems into intermediate steps, allowing more computational resources to be dedicated to problems that involve multiple reasoning stages.
    \item Secondly, chain of thought offers a transparent view into the model’s decision-making process, indicating how a particular solution may have been reached and enabling debugging of the reasoning path when errors occur (though fully understanding all computations behind an answer remains unresolved).
    \item Thirdly, chain-of-thought reasoning is applicable to tasks like math word problems, commonsense reasoning, and symbolic manipulation, and potentially can be extended (at least in theory) to any task solvable by humans through language.
    \item Lastly, chain-of-thought reasoning can be easily prompted in sufficiently large pretrained language models by simply incorporating chain-of-thought example sequences into few-shot prompt exemplars.
\end{enumerate}

In our experiments, we will demonstrate the effectiveness of chain-of-thought prompting for arithmetic reasoning (\cref{sec:arithmetic-reasoning}), commonsense reasoning (\cref{sec:commonsense-reasoning}), and symbolic reasoning (\cref{sec:symbolic-reasoning}).

\section{Arithmetic Reasoning}\label{sec:arithmetic-reasoning}
We start by examining math word problems as illustrated in \cref{fig:pull-figure}, which test the arithmetic reasoning capabilities of language models. While these problems are generally straightforward for humans, language models frequently find arithmetic reasoning challenging \citep[][\textit{inter alia}]{hendrycks2021measuring,patel-etal-2021-nlp}. Remarkably, when combined with a 540B parameter language model, chain-of-thought prompting matches the performance of task-specific finetuned models on several datasets and even sets a new state-of-the-art on the difficult GSM8K benchmark \citep{cobbe2021training}.

\subsection{Experimental Setup}

We evaluate chain-of-thought prompting across different language models on a variety of benchmarks.

\textbf{Benchmarks.}
We use five math word problem datasets:
\textbf{(1)} the \textbf{GSM8K} benchmark of math word problems \citep{cobbe2021training}, 
\textbf{(2)} the \textbf{SVAMP} dataset featuring math word problems with diverse structures \citep{patel-etal-2021-nlp},
\textbf{(3)} the \textbf{ASDiv} dataset containing a variety of math word problems \citep{miao-etal-2020-diverse},
\textbf{(4)} the \textbf{AQuA} dataset with algebraic word problems, and
\textbf{(5)} the \textbf{MAWPS} benchmark \citep{koncel-kedziorski-etal-2016-mawps}. Sample problems are provided in Appendix \cref{tab:math-datasets}.

\textbf{Standard prompting.} As a baseline, we employ standard few-shot prompting as introduced by \citet{brown2020language}, where a language model is given in-context examples of input-output pairs before generating an answer for a test query. The exemplars are presented as questions and answers, and the model directly produces the answer, as depicted in \cref{fig:pull-figure} (left).

\textbf{Chain-of-thought prompting.} Our method involves enhancing each exemplar in few-shot prompting by including a chain of thought corresponding to the answer, as depicted in \cref{fig:pull-figure} (right). Since most datasets provide only an evaluation split, we manually created a set of eight few-shot exemplars accompanied by chains of thought for prompting—one such chain of thought exemplar is shown in \cref{fig:pull-figure} (right), and the complete set of exemplars is listed in Appendix \cref{tab:appendix-mwp-prompts}. (These exemplars were not subjected to prompt engineering; robustness is analyzed in \cref{subsec:robustness} and \cref{subsec:faq-prompt-engineering}.) To examine whether this style of chain-of-thought prompting effectively triggers successful reasoning across various math word problems, we applied this single set of eight chain-of-thought exemplars to all benchmarks except AQuA, which uses a multiple-choice format rather than free response. For AQuA, we employed four exemplars and solutions drawn from the training set, detailed in Appendix \cref{tab:appendix-aqua-prompt}.

\textbf{Language models.} We assess five large language models. The first is \textbf{GPT-3} \citep{brown2020language}, for which we utilize text-ada-001, text-babbage-001, text-curie-001, and text-davinci-002, likely corresponding to InstructGPT models with 350M, 1.3B, 6.7B, and 175B parameters respectively \citep{ouyang2022training}. The second model family is \textbf{\lamda{}} \citep{thoppilan2022lamda}, with models sized 422M, 2B, 8B, 68B, and 137B parameters. The third is \textbf{\palm{}}, featuring models of 8B, 62B, and 540B parameters. The fourth is \textbf{UL2 20B} \citep{tay2022unifying}, and the fifth is \textbf{Codex} \citep[][code-davinci-002 in the OpenAI API]{chen2021evaluating}. We generate outputs from the models using greedy decoding (although subsequent work demonstrates that chain-of-thought prompting performance can be enhanced by taking the majority vote over multiple sampled generations \citep{wang2022self}). For \lamda{}, results are averaged over five random seeds, each corresponding to a different randomized order of exemplars. Since \lamda{} experiments exhibited low variance across seeds, to reduce computational cost, we report results using a single exemplar ordering for all other models.

\subsection{Results}
The most notable outcomes of chain-of-thought prompting are highlighted in \cref{fig:main-math}, with comprehensive experimental results for each model collection, model size, and benchmark provided in \cref{tab:all-lm-math} in the Appendix. There are three main conclusions. First, \cref{fig:main-math} demonstrates that chain-of-thought prompting is an emergent property related to model scale \citep{wei2022emergent}. Specifically, chain-of-thought prompting does not enhance performance for smaller models and only starts to improve results when applied to models with approximately 100B parameters. Our qualitative analysis showed that smaller models generated coherent but illogical chains of thought, resulting in worse performance compared to standard prompting.

\input{main_fables/main-math}
Second, chain-of-thought prompting leads to greater performance improvements on more complex problems. For example, in GSM8K (the dataset with the lowest baseline accuracy), performance more than doubled for the largest GPT and \palm{} models. Conversely, for SingleOp, the simplest subset of MAWPS requiring only a single step to solve, performance gains were minimal or even negative (see Appendix \cref{tab:mawps-subsets}).

Third, chain-of-thought prompting using GPT-3 175B and \palm{} 540B achieves results competitive with previous state-of-the-art methods, which typically involve fine-tuning a task-specific model on labeled training data. 
\cref{fig:main-math} illustrates how \palm{} 540B leverages chain-of-thought prompting to set new state-of-the-art performance on GSM8K, SVAMP, and MAWPS (although note that standard prompting already surpassed the prior best on SVAMP). On the remaining two datasets, AQuA and ASDiv, \palm{} with chain-of-thought prompting attains results within 2\% of the state of the art (see Appendix \cref{tab:all-lm-math}).

To gain a deeper understanding of why chain-of-thought prompting is effective, we conducted a manual inspection of the chains of thought generated by \lamda{} 137B for GSM8K. Among 50 randomly selected instances where the model produced the correct final answer, every generated chain of thought was logically and mathematically sound except for two cases where the correct answer was coincidentally obtained (refer to \cref{subsec:correct-chain-of-thought} and \cref{tab:appendix-gsm8k-correct-examples} for samples of correct chains of thought generated by the model). Additionally, we reviewed 50 random samples where the model provided an incorrect answer. Our analysis revealed that 46\% of these chains were nearly correct, aside from minor errors such as calculator mistakes, symbol mapping issues, or a single reasoning step missing, while the remaining 54\% contained significant errors related to semantic understanding or coherence (see \cref{subsec:incorrect-chain-of-thought-analysis}). 

To shed some light on why scaling enhances chain-of-thought reasoning capabilities, we performed a comparable error analysis on \palm{} 62B and examined whether these errors were resolved when scaling up to \palm{} 540B. The findings indicate that increasing \palm{} to 540B mitigates a substantial number of one-step missing errors and semantic understanding mistakes present in the 62B model (see \cref{subsec:why-scale-helps}).

\subsection{Ablation Study}\label{subsec:math-ablation}

The improvements observed from chain-of-thought prompting naturally prompt the question of whether similar performance gains can be achieved using alternative types of prompting. 

\cref{fig:bar_ablation} presents an ablation study involving three variations of chain-of-thought prompting, described below.

\textbf{Equation only.} One potential reason chain-of-thought prompting is beneficial is that it generates the mathematical equation to be solved. To test this, we examine a variant where the model is instructed to output only the mathematical equation before providing the final answer. 

As shown in \cref{fig:bar_ablation}, equation-only prompting offers limited improvement on GSM8K, suggesting that the semantics of GSM8K questions are too complex to be directly converted into equations without the intermediate natural language reasoning steps characteristic of chain-of-thought prompting. However, for datasets involving one-step or two-step problems, equation-only prompting does enhance performance, as the equations can be readily derived from the questions (see Appendix \cref{tab:ablations-arithmetic}).

\input{main_fables/ablation-bar}
\textbf{Variable compute only.} One hypothesis is that the chain of thought enables the model to allocate more computational effort (i.e., produce more intermediate tokens) on more challenging problems. To separate the impact of variable computation from the reasoning process in the chain of thought, we experiment with a setting where the model is instructed to produce only a sequence of dots ($\ldots$) matching the number of characters in the equation necessary to solve the problem. This variation performs roughly the same as the baseline, indicating that variable computation alone does not account for the effectiveness of chain-of-thought prompting, and that articulating intermediate steps in natural language offers additional benefits.

\textbf{Chain of thought after answer.} Another possible advantage of chain-of-thought prompting might be that it helps the model better retrieve relevant knowledge learned during pretraining. To examine this, we test a different setup where the chain-of-thought prompt is provided only after the model generates the answer, thereby isolating whether the model relies on the generated chain of thought to produce the final response. This variant yields performance similar to the baseline, suggesting that the sequential reasoning represented by the chain of thought contributes value beyond merely activating memorized knowledge.

\input{main_fables/main-robustness}
\subsection{Robustness of Chain of Thought}\label{subsec:robustness}

Prompting methods are highly sensitive to the choice of exemplars—for example, changing the order of few-shot exemplars can cause GPT-3's accuracy on SST-2 to vary dramatically from near chance (54.3\%) to nearly state-of-the-art levels (93.4\%) \citep{zhao2021calibrate}. In this final subsection, we investigate the robustness of chain of thought prompts authored by different annotators. Beyond the previously reported results that used chains of thought created by Annotator A, two other co-authors of this study (Annotators B and C) independently authored chains of thought for the same few-shot examples (see \cref{sec:appendix-alternate-annotators}). Annotator A also authored a more concise version of the chain of thought, modeled after the style of solutions found in \citet{cobbe2021training}.\footnote{For example, while the original chain of thought uses several brief sentences (\textit{``There were originally 9 computers. For each of 4 days, 5 more computers were added. So 5 * 4 = 20 computers were added. 9 + 20 is 29.''}), the concise chain reads \textit{``5 * 4 = 20 new computers were added. So there are 9 + 20 = 29 new computers in the server room now.''}}

\cref{fig:bar-robustness-analysis} presents these findings for \lamda{} 137B on GSM8K and MAWPS (additional ablation results for other datasets can be found in Appendix \cref{tab:ablations-arithmetic} / \cref{tab:ablations-commonsense-symbolic}). Although there is variability across different chain of thought annotations, as expected when employing exemplar-based prompting \citep{le-scao-rush-2021-many,reynolds2021prompt,zhao2021calibrate}, every set of chain of thought prompts significantly outperforms the standard baseline. This indicates that the effective use of chain of thought is not reliant on any specific linguistic style.

To verify that effective chain-of-thought prompting holds for other exemplar sets, we also conducted experiments using three sets of eight exemplars randomly drawn from the GSM8K training set, which is an independent source (the examples in this dataset already contained reasoning steps akin to a chain of thought).\footnote{We select examples with $\leq 60$ tokens to fit within our input context window, and also restrict the examples to $\leq 2$ steps to allow a fair comparison with the eight exemplars we composed.} \cref{fig:bar-robustness-analysis} demonstrates that these prompts achieved similar performance to our manually crafted exemplars, again substantially surpassing standard prompting.

Beyond robustness to annotator variation, independently-written chains of thought, different exemplars, and diverse language models, we also observe that chain-of-thought prompting for arithmetic reasoning is resilient to changes in exemplar order and variations in the number of exemplars used (see \cref{subsec:faq-prompt-engineering}).

\section{Commonsense Reasoning}\label{sec:commonsense-reasoning}

While chain of thought is especially well-suited for math word problems, its language-based nature enables it to be applied to a wide range of commonsense reasoning tasks, which require reasoning about physical and human interactions based on assumed general background knowledge. Commonsense reasoning is essential for engaging with the world and remains beyond the capabilities of current natural language understanding systems \citep{talmor2022commonsenseqa}.

\textbf{Benchmarks.}  
We examine five datasets that encompass a wide variety of commonsense reasoning challenges. The well-known \textbf{CSQA} \citep{talmor-etal-2019-commonsenseqa} presents commonsense questions about the world, involving intricate semantics that often necessitate background knowledge.  
\textbf{StrategyQA} \citep{geva-etal-2021-aristotle} demands that models deduce a multi-step reasoning strategy to answer questions. We select two specialized evaluation sets from the BIG-bench collection \citep{bigbench}: \textbf{Date} Understanding, which requires inferring a date based on a given context, and \textbf{Sports} Understanding, which focuses on judging whether a sports-related sentence is plausible or not. Lastly, the \textbf{SayCan} dataset \citep{ahn2022can} involves translating natural language instructions into a sequence of discrete robot actions.  
\cref{fig:dataset-examples} presents examples with chain of thought annotations across all datasets.

\textbf{Prompts.}  
We adopt the same experimental methodology as described in the previous section. For CSQA and StrategyQA, we randomly picked samples from the training sets and manually crafted chains of thought to serve as few-shot exemplars. Since the two BIG-bench tasks lack training sets, we selected the first ten examples from their evaluation sets as few-shot exemplars and report results on the remaining evaluation examples. For SayCan, we utilized six examples from the training set used by \citet{ahn2022can} and also manually generated chains of thought.

\textbf{Results.}
\cref{fig:commonsense-results} presents these findings for \palm{} (complete results for \lamda{}, GPT-3, and various model sizes are provided in \cref{tab:all-lm-commonsense}).
Across all tasks, increasing model size enhanced the effectiveness of standard prompting; employing chain-of-thought prompting yielded additional improvements, with the greatest gains observed in \palm{} 540B. Using chain-of-thought prompting, \palm{} 540B demonstrated strong performance compared to baselines, surpassing the previous state of the art on StrategyQA (75.6\% versus 69.4\%) and exceeding the accuracy of an unaided sports fan in sports comprehension (95.4\% versus 84\%).
These outcomes indicate that chain-of-thought prompting can also boost performance on tasks involving diverse commonsense reasoning skills (although the improvement was minimal on CSQA).

\input{main_fables/main-commonsense}

\input{main_fables/main-symbolic}
\section{Symbolic Reasoning}\label{sec:symbolic-reasoning}
Our final set of experiments focuses on symbolic reasoning, which humans find straightforward but may pose difficulties for language models. We demonstrate that chain-of-thought prompting not only allows language models to tackle symbolic reasoning tasks that are hard to solve with standard prompting but also supports length generalization to input sequences during inference that are longer than those presented in few-shot examples.

\paragraph{Tasks.}
We employ the following two simplified tasks.

\begin{itemize}[leftmargin=*,topsep=0pt]
    \itemsep0em \item \textbf{Last letter concatenation.} This task requires the model to join the last letters of words in a given name (e.g., \textit{``Amy Brown''} $\rightarrow$ \textit{``yn''}). 
    It is a more difficult variant of first letter concatenation, a task that language models can already solve without chain of thought.\footnote{We evaluated 10 common names using GPT-3 \texttt{davinci} and it succeeded on all but one.} Full names are generated by randomly combining names from the top one-thousand most common first and last names based on name census data ({\small{\url{https://namecensus.com/}}}).
    \item \textbf{Coin flip.} This task involves determining whether a coin remains heads up after a sequence of people either flip or do not flip the coin (e.g., \textit{``A coin is heads up. Phoebe flips the coin. Osvaldo does not flip the coin. Is the coin still heads up?''} $\rightarrow$ \textit{``no''}).
\end{itemize}

Given that the design of these symbolic reasoning tasks is clearly specified, for each task we examine an \textit{in-domain} test set where the examples have the same number of steps as the training or few-shot examples, as well as an \textit{out-of-domain} (OOD) test set, where the evaluation examples involve more steps than those in the exemplars. For the last letter concatenation task, the model is exposed only to exemplars of names containing two words, and then it performs last letter concatenation on names with 3 and 4 words.\footnote{For names longer than two words, we concatenate several first and last names together.} A similar approach is taken for the number of possible flips in the coin flip task. Our experimental setup uses the same techniques and models as described in the previous two sections. We once again manually create chains of thought for the few-shot examples for each task, which are presented in \cref{fig:dataset-examples}.

\paragraph{Results.} 
The outcomes of these in-domain and OOD evaluations are displayed in \cref{fig:symbolic-main} for \palm{}, with \lamda{} results included in Appendix \cref{tab:all-lm-symbolic}. Using \palm{} 540B, chain-of-thought prompting yields nearly 100\% success rates (noting that standard prompting already solves the coin flip task with \palm{} 540B, though this is not the case for \lamda{} 137B). It is important to highlight that these in-domain tests are ``toy tasks'' in the sense that the perfect solution frameworks are already provided by the chains of thought in the few-shot exemplars; the model simply needs to replicate the same steps with the new symbols during testing. Nevertheless, smaller models still fail— the capability to perform abstract manipulations on previously unseen symbols for these three tasks only emerges at around the 100B parameter scale.

Regarding the OOD evaluations, standard prompting does not succeed for either task. However, with chain-of-thought prompting, language models display upward scaling trends (although their performance remains lower than in the in-domain context). Therefore, chain-of-thought prompting enables length generalization beyond the previously encountered chains of thought for sufficiently large language models.

\section{Discussion}\label{sec:discussion}

We have investigated chain-of-thought prompting as a straightforward approach to trigger multi-step reasoning capabilities in large language models. Initially, we observed that chain-of-thought prompting significantly enhances performance on arithmetic reasoning, producing gains that are substantially greater than ablation studies and remain consistent across various annotators, exemplars, and language models (\cref{sec:arithmetic-reasoning}). Subsequently, experiments on commonsense reasoning highlighted that the linguistic nature of chain-of-thought reasoning makes it broadly applicable (\cref{sec:commonsense-reasoning}). Lastly, we demonstrated that for symbolic reasoning, chain-of-thought prompting supports OOD generalization to longer sequences (\cref{sec:symbolic-reasoning}). Across all experiments, chain-of-thought reasoning was elicited simply by prompting an off-the-shelf language model, without any finetuning performed during this study.

The appearance of chain-of-thought reasoning as a function of model scale has been a notable trend \citep{wei2022emergent}. For many reasoning tasks where standard prompting yields a flat scaling curve, chain-of-thought prompting induces substantially rising scaling curves. Chain-of-thought prompting seems to broaden the set of tasks that large language models can solve effectively—in other words, our findings emphasize that standard prompting represents only a lower bound on the abilities of large language models. This insight likely prompts more questions than it resolves—for example, to what extent can reasoning capabilities improve with further increases in model scale? What other prompting strategies might expand the repertoire of tasks that language models can tackle?

Concerning limitations, we first clarify that although chain-of-thought mimics human reasoning processes, it does not confirm whether the neural network is genuinely "reasoning," which remains an open question. Second, while the manual effort to augment exemplars with chains of thought is modest in the few-shot setting, such annotation costs might be prohibitive for finetuning (though this might be mitigated through synthetic data generation or zero-shot generalization). \orange{Third, there is no assurance that the reasoning paths are correct, which can result in both correct and incorrect answers; enhancing the factual accuracy of language model outputs is an ongoing area for future research \citep[][\textit{inter alia}]{rashkin2021measuring,ye2022unreliability,wiegreffe2021reframing}.} Lastly, since chain-of-thought reasoning emerges only at large model scales, deploying it in practical applications can be expensive; future work might investigate methods to induce reasoning in smaller models.

\section{Related Work} This study draws inspiration from multiple research areas, which we discuss more comprehensively in the extended related work section (\cref{sec:extended-related-work}). Here, we focus on two key directions and the corresponding literature that are arguably the most pertinent.

The first relevant line of research involves utilizing intermediate steps to tackle reasoning tasks. \cite{ling-etal-2017-program} introduce the concept of employing natural language rationales to address math word problems by breaking down the solution into a sequence of intermediate steps. Their approach stands in stark contrast to prior work that relies on formal languages for reasoning \citep{roy-2016-reasoning, chiang-chen-2019-semantically, amini-etal-2019-mathqa, chen2019neural}. Building on \citet{ling-etal-2017-program}, \citet{cobbe2021training} expand the dataset and utilize it to finetune a pretrained language model instead of training one from the ground up. In program synthesis, \cite{nye2021show} exploit language models to predict the final outputs of Python programs by first forecasting intermediate computational results line-by-line, demonstrating that this stepwise prediction approach outperforms directly predicting the final outputs.

Additionally, this work is closely connected to a significant body of recent research on prompting techniques. Since the introduction of few-shot prompting by \citet{brown2020language}, various general methods have been proposed to enhance the prompting capabilities of models, such as automatically generating prompts \citep{lester-etal-2021-power} or providing models with task descriptions through instructions \citep{wei2021finetuned,sanh2021multitask,ouyang2022training}. While these methods primarily improve or modify the input portion of the prompt (e.g., instructions prepended to inputs), our approach takes a complementary direction by enriching the outputs of language models with a chain of thought.

\section{Conclusions} We have investigated chain-of-thought prompting as a straightforward and widely applicable technique for boosting reasoning capabilities in language models. Our experiments across arithmetic, symbolic, and commonsense reasoning tasks reveal that chain-of-thought reasoning emerges as a property of model scale, enabling sufficiently large language models to succeed in reasoning tasks that otherwise show flat performance improvements with scale. Expanding the scope of reasoning tasks that language models can handle may encourage further research on language-centric methods for reasoning.

\end{document}